\subsection{The residual order-two dyadic branch}

\begin{lemma}[B-RD47: the two quadratic fields of a two-isogeny]
\label{lem:B-RD47}
\status{Proved; complete independent parent and monodromy proof and exact-code audits passed.}
Let $K$ be complete, with algebraically closed residue field of
characteristic two and $v_K(2)=1$. Suppose $E/K$ is potentially
supersingular, and both $E$ and its quotient by a rational
order-two subgroup have nontrivial residual image $C_2$.
Write
\[
 E:\quad y^2=x(x^2+a x+b),\qquad D=a^2-4b.
\]
The two quadratic two-torsion fields are $K(\sqrt D)$ and
$K(\sqrt b)$. The special-inertia image has order divisible by
three if and only if $v_K(D)$ and $v_K(b)$ have different parities.
\end{lemma}
\begin{proof}
Completing the square and translating the rational two-torsion
point gives the displayed equation, with $bD\ne0$.
The normalized quotient is
\[
 E_1:\quad Y^2=X(X^2-2aX+D),\qquad
 X=x+a+b/x,\quad Y=y(1-b/x^2).                         \tag{B-RD47.1}
\]
Direct substitution verifies the equation and
$dX/Y=dx/y$. The map has degree two, sends the origin and
$(0,0)$ to the origin, and hence is the quotient isogeny.
Its source two-torsion field is $K(\sqrt D)$; the target
quadratic discriminant is $16b$, giving $K(\sqrt b)$.
The hypotheses say that both $D$ and $b$ are nonsquares.

We shall use
\[
                         1+4\mathcal O_K
                              \subset K^{\times2}.     \tag{B-RD47.2}
\]
Indeed the equation $z^2+z=u$ has a residue solution for every
$u\in\mathcal O_K$, and that solution lifts by Hensel's lemma
since $2z+1$ is a unit. Then $(1+2z)^2=1+4u$.

A quadratic twist replaces $(a,b,D)$ by $(ta,t^2b,t^2D)$.
It preserves the two quadratic fields, their valuation
parities, and the discriminant valuation modulo three.
Choose such a twist to normalize
$\beta=v_K(b)\in\{0,1\}$, and write $\alpha=v_K(a)$,
allowing $\alpha=+\infty$ when $a=0$.
The exact invariants are
\[
 \Delta=16b^2D,\qquad c_4=16(a^2-3b),\qquad
 v_K(j)=8+3v_K(a^2-3b)-2\beta-v_K(D).                 \tag{B-RD47.3}
\]

If $\beta=0$ and $\alpha\le0$, then
$D=a^2(1-4b/a^2)$ is a square by (B-RD47.2), contrary to the
hypothesis. Thus $\alpha\ge1$. For $\alpha\ge2$ one has
$v_K(D)=2$ and $v_K(a^2-3b)=0$, so
$(v_K(\Delta),v_K(j))=(6,6)$.
For $\alpha=1$, put $a=2c$ with $c$ a unit and
$r=v_K(c^2-b)\ge0$.
If $r=0$, the same pair is $(6,6)$.
If $r\ge2$, the identity
$b=c^2(1-(c^2-b)/c^2)$ and (B-RD47.2) make $b$ a square,
again contrary to the hypothesis. Hence the only other case
is $r=1$, when $v_K(D)=3$, $v_K(a^2-3b)=0$, and the pair
is $(7,5)$.

If $\beta=1$ and $\alpha\le0$, the same factorization
$D=a^2(1-4b/a^2)$ makes $D$ a square. Thus $\alpha\ge1$.
For $\alpha=1$, its terms have valuations two and three,
so $v_K(D)=2$; also $v_K(a^2-3b)=1$. The invariant pair
is $(8,7)$. For $\alpha\ge2$, one instead has
$v_K(D)=3$, $v_K(a^2-3b)=1$, giving $(9,6)$.
We have exhausted all possibilities under the hypotheses:
\[
\begin{array}{c|c|c|c|c|c}
 \beta&\text{case}&v_K(D)&v_K(\Delta)&v_K(j)
       &3\mid |H|\\\hline
 0&\alpha\ge2\text{ or }(\alpha,r)=(1,0)&2&6&6&\text{no}\\
 0&(\alpha,r)=(1,1)&3&7&5&\text{yes}\\
 1&\alpha=1&2&8&7&\text{yes}\\
 1&\alpha\ge2&3&9&6&\text{no}
\end{array}                                                   \tag{B-RD47.4}
\]
The last column is Lemma~\ref{lem:B-RD46-discriminant}.
It is precisely the assertion about the two valuation parities.
\end{proof}

\begin{proposition}[B-RD47: order-two residual inertia is excluded]
\label{prop:B-RD47}
\status{Proved; complete independent parent and monodromy proof and exact-code audits passed.}
Let $K$ be a complete discretely valued field of mixed characteristic
$(0,2)$, with perfect residue field and $v_K(2)=1$. Suppose two
potentially supersingular elliptic curves have isomorphic
eight-torsion and their common residual inertia image is $C_2$.
Their special-inertia characters agree on elements of odd
order in a common finite special-inertia quotient.

Consequently every hypothetical height-two, ramification-one
sixteen-torsion ambiguity left by
Proposition~\ref{prop:B-RD44} has common residual inertia image
exactly $C_3$. No existence or nonexistence in that final branch
is asserted here.
\end{proposition}
\begin{proof}
Pass to the completed maximal unramified extension, preserving
the given isomorphism, the valuation, and the special character.
The unique nonzero invariant vector in the common residual
$C_2$-module identifies rational order-two subgroups
$C\subset E[2]$ and $C'\subset E'[2]$ under the given
isomorphism. Put $E_1=E/C$ and $E'_1=E'/C'$.
There is an exact scheme-theoretic construction
\[
 T=\{P\in E[8]:[4]P\in C\},\qquad
                         T/C\simeq E_1[4].             \tag{B-RD47.5}
\]
Indeed the preimage of $E_1[4]$ under the quotient isogeny is
exactly $T$; if $[4]P\in C$ then $[8]P=0$.
The given eight-torsion isomorphism therefore induces
$E_1[4]\simeq E'_1[4]$. In particular their two-torsion
representations agree.

Each quotient has a rational two-torsion point, namely the
dual-isogeny kernel, so its residual image is either trivial
or $C_2$. If it is trivial, both quotient curves have rational
full two-torsion, and Proposition~\ref{prop:B-RD46} applies to
their four-torsion isomorphism. Isogeny preserves the rational
special-inertia character, proving the assertion in this case.

Otherwise both quotient residual images are $C_2$. The
original two-torsion representations identify the nontrivial
quadratic fields $K(\sqrt D)$ and $K(\sqrt {D'})$, while the
quotient two-torsion representations identify
$K(\sqrt b)$ and $K(\sqrt {b'})$ in the equations of
Lemma~\ref{lem:B-RD47}. Equality of quadratic extensions
is equality of their nonzero Kummer classes, so it preserves
the valuation parity of each class. That lemma gives equality
of the cubic special-inertia conditions, hence equality of the
requested characters by Lemma~\ref{lem:B-RD46-discriminant}.

For the final statement, use the exact elliptic realizations
and mixed exclusion of Proposition~\ref{prop:B-RD44}, restrict
the sixteen-torsion isomorphism to eight-torsion, and combine
this proposition with the trivial-residual exclusion of
Proposition~\ref{prop:B-RD46}.
\end{proof}

\paragraph{Exact verification.}
The retained program
\path{code/verify_p2_residual_c2_gate.py}
verifies the quotient equation and normalized differential,
the two discriminants and $j$ formulas, and every row and
excluded-square case of the valuation table. The proof of
the finite-group-scheme quotient in (B-RD47.5) is given
above. No search of the remaining residual $C_3$ branch is
represented by that certificate. Both statements have passed complete independent audits.

The author certificate passed in $0.20210709299863083$ seconds.

\paragraph{Independent audits.}
The parent independently checked the full proof and exact code,
rerunning the certificate in $0.047221027001796756$ seconds
at 2026-11-19 14:16:15.762965 UTC. The monodromy agent completed
its own full proof and code audit, with an author-preserving
rerun of $0.15887224199832417$ seconds. Both checked the
finite-group-scheme quotient and the entire valuation table.

\paragraph{Subsequent closure.}
The remaining residual branches are settled in
Theorem~\ref{thm:B-RD48}, which gives the exact value
$\nu_{\min}(2,1,2)=3$. The assertions and finite certificate
in this subsection retain their stated narrower scope.
